Los experimentos fueron ejecutados en un entorno bajo el sistema operativo Ubuntu 22.04.5 LTS (WSL2 o Nativo).
Esta cuenta con un procesador AMD Ryzen 5300U con AMD Radeon Graphics x8, 8 GB de RAM y disco sólido de 256 GB.
Todo el código fuente fue desarrollado en \texttt{C++}, siendo compilado con \texttt{g++} apoyándose en un
Makefile que automatizó los procesos de limpieza, construcción, generación de casos de prueba y ejecución
de mediciones. Una decisión técnica de gran relevancia fue aislar el consumo de memoria mediante llamadas
al sistema operativo (\texttt{fork} y \texttt{wait4}), lo cual nos permitió utilizar \texttt{getrusage()}
para obtener el pico de memoria (\texttt{Maximum Resident Set Size}) verdaderamente originado por cada algoritmo,
 sin verse afectado por la memoria dinámica alojada durante la lectura de los datos de entrada o ejecuciones anteriores.

Para los datasets de evaluación, se implementó el script \texttt{testcases\_generator.py} según las instrucciones del
enunciado, capaz de producir escenarios variados. Se generaron casos pequeños (donde $N < 15$) utilizados para contrastar
Fuerza Bruta contra Programación Dinámica (DP), y casos medianos a grandes (con $N$ rondando los $100$ y hasta cientos de
capítulos en total) pensados para exponer las diferencias de tiempo, memoria y porcentaje de optimalidad entre DP y los
enfoques Greedy.
